\section{Introduction}
The computational demands of large language models (LLMs) continue to grow rapidly, driving training and inference costs to millions of dollars. High-performance GPU kernels are critical to reducing these costs by determining the throughput and efficiency of modern deep learning workloads. However, developing efficient GPU kernels remains heavily reliant on manual tuning by domain experts, a process that is increasingly expensive and time-consuming as GPU architectures evolve and diversify. These challenges motivate automated GPU kernel generation as a critical research problem.

Recent advances in LLMs have demonstrated strong capabilities in general-purpose code generation (e.g., Python, C++), making them promising candidates for automating GPU kernel generation. However, GPU kernel generation is fundamentally different from general-purpose programming. It is not merely a code synthesis task, but a tightly constrained optimization problem over a vast, hardware-dependent design space. Kernel code must satisfy strict syntactic and semantic correctness requirements while simultaneously achieving high performance. As a result, directly applying LLMs without additional information often produces kernels that are invalid or far from optimal, motivating specialized methodological designs to bridge this gap.

Accordingly, existing approaches for LLM-based GPU kernel generation can be broadly categorized into two classes. The first class specializes LLMs through fine-tuning or reinforcement learning (RL) to adapt them to the kernel generation domain. The second class builds agent-based systems on top of general-purpose pretrained LLMs, leveraging iterative refinement to improve generated kernels without additional model training.

Despite their differences, both classes of approaches are fundamentally constrained by the extreme complexity of the GPU kernel optimization space. Kernel performance is governed by intricate interactions among thread scheduling, memory hierarchy utilization, synchronization mechanisms, and hardware-specific features. Even minor code modifications can result in orders-of-magnitude performance differences. Prior work shows that even a small neural network subgraph can expose an optimization space of up to $10^9$
 configurations on GPUs \fktocheck{(Zhai et al., 2024)}. Without expert knowledge to systematically constrain and navigate this vast space, automated methods struggle to reliably converge to high-performance solutions.

Under this challenge, RL-based approaches face significant limitations. High-quality GPU kernel datasets are scarce, and kernel code must satisfy strict syntactic, semantic, and hardware-specific constraints. Consequently, RL-generated kernels frequently fail to compile or produce incorrect results, leading to low overall effectiveness. For example, AutoTriton reports an average correctness rate below 40\%, highlighting the difficulty of directly applying RL to GPU kernel generation.

Agent-based approaches built on general-purpose LLMs offer a more practical alternative by leveraging pretrained models' capabilities without additional training. Representative systems such as GEAK employ multi-turn iterative refinement to improve kernel correctness and robustness, resulting in better engineering feasibility and scalability. However, most existing agent systems still lack domain-specific kernel optimization knowledge and rely primarily on coarse-grained performance feedback (e.g., execution time). As a result, agents often perform unguided trial-and-error exploration, which severely limits achievable performance gains; for instance, GEAK exhibits marginal improvement after multiple optimization rounds.

These limitations suggest that fully unlocking the potential of LLMs for GPU kernel generation requires integrating expert-level kernel optimization knowledge to systematically guide the exploration process. Rather than relying on unguided trial-and-error or end-to-end learning, effective systems must impose structure on the optimization space that reflects expert reasoning and hardware-aware constraints.

 To this end, we propose XX, an expert-guided agent framework for automatic generation of high-performance GPU kernels. XX incorporates domain knowledge into the agent decision process through a staged optimization strategy, enabling structured exploration of the kernel optimization space while preserving the flexibility of general-purpose LLMs.

 Specifically, XX employs expert-guided staged optimization, decomposing the kernel generation workflow into two hierarchical levels: algorithmic structure design and hardware-specific tuning. The algorithmic structure design stage establishes a strong performance upper bound by combining multi-seed search with algorithmic refinement techniques. The hardware-specific tuning stage then systematically realizes this potential through three sequential sub-stages with explicit objectives: parallel mapping, tensor tiling, and memory optimization. By decomposing the complex optimization problem into constrained sub-problems aligned with expert knowledge, XX effectively guides LLM-based agents toward high-performance kernel implementations.

 To support this workflow, we propose a stage-aware multi-agent collaboration mechanism consisting of a code agent, a profiling agent, and a debugging agent. Within each stage, the agents coordinate around clearly defined objectives via structured information exchange, including bottleneck identification and optimization proposals. Across stages, a selective context propagation strategy retains only finalized optimization decisions while discarding intermediate exploratory artifacts. This design reduces cross-stage objective interference and enables consistent, cumulative performance improvements over the course of the optimization process.

We systematically evaluate our approach on KernelBench and real-world workloads. Experimental results show that our approach achieves 100\% correctness while delivering up to $2.13\times$ speedup over PyTorch implementations, demonstrating the effectiveness of integrating expert knowledge into the kernel generation agent.

The key contributions of this paper are as follows:
\begin{itemize}
\item We propose an expert-guided agent framework that decomposes kernel optimization into staged sub-problems (algorithmic design and hardware-specific tuning), effectively constraining the vast GPU optimization space.

\item We design a stage-aware multi-agent collaboration mechanism with selective context propagation that achieves cumulative, interference-free improvements across optimization stages.

\item We achieve 100\% correctness and up to $2.13\times$ speedup over PyTorch on KernelBench and real-world workloads, demonstrating effective integration of expert knowledge into LLM-based kernel generation.
\end{itemize}
